%%%%%%%%%%%%%%%%%%%%%%%%%%%%%%%%%%%%%%%%%%%%%%%%%%%%%%%%%%%%%%%%%%%%%%%%%%%%%%%%
% kitchen-sink.tex — the parser/WYSIWYG conformance fixture.
%
% Unlike heavy-tail-paper.tex (a realistic paper), this file is a deliberate
% torture test: every construct here is one the WYSIWYG view has to survive
% without dropping, reordering, or rewriting. Anything that regresses shows up
% as a failing assertion in test/parser.kitchen-sink.test.ts or as drift in
% test/roundtrip.stability.test.ts.
%
% Compiles with pdflatex (three passes for cleveref/toc).
%%%%%%%%%%%%%%%%%%%%%%%%%%%%%%%%%%%%%%%%%%%%%%%%%%%%%%%%%%%%%%%%%%%%%%%%%%%%%%%%

\documentclass[11pt,twoside]{article}

\usepackage[utf8]{inputenc}
\usepackage[T1]{fontenc}
\usepackage[margin=1in]{geometry}
\usepackage{amsmath, amssymb, amsthm, mathtools}
\usepackage{booktabs}
\usepackage{multirow}
\usepackage{array}
\usepackage{longtable}
\usepackage{graphicx}
\usepackage{subcaption}
\usepackage{enumitem}
\usepackage{xcolor}
\usepackage{listings}
\usepackage{natbib}
\usepackage{hyperref}
\usepackage[capitalise]{cleveref}

\newtheorem{theorem}{Theorem}[section]
\newtheorem{lemma}[theorem]{Lemma}
\newtheorem{definition}[theorem]{Definition}

\DeclareMathOperator{\Tr}{Tr}
\DeclareMathOperator*{\argmin}{arg\,min}
\newcommand{\R}{\mathbb{R}}
\newcommand{\norm}[1]{\left\lVert#1\right\rVert}

\title{A Conformance Suite for \LaTeX{} Round-Tripping}
\author{Dana Ortiz\thanks{Thanks to the reviewers.} \and Ren\'{e} M\"{u}ller}
\date{\today}

\begin{document}
\maketitle

\begin{abstract}
This document exercises constructs that a WYSIWYG editor must round-trip
byte-for-byte: tabular material, nested lists, multi-line math with per-line
labels, verbatim text, accents, and escaped characters.
\end{abstract}

\tableofcontents

\section{Escapes and Text-Mode Typography}
\label{sec:escapes}

Reserved characters must survive unchanged: 100\% of \$5 costs \#1 in the
Smith \& Jones catalogue, filed under \texttt{report\_final}, with braces
\{like this\} and a literal backslash \textbackslash{} plus a tilde
\textasciitilde{} and caret \textasciicircum{} for good measure.

Ligature-sensitive punctuation: an em-dash---like so---and an en-dash 3--7,
with ``curly quotes'' and `single quotes' around them. A non-breaking
space appears in \cref{sec:escapes}~and Table~\ref{tab:results}.

Accented text set with commands: Ren\'{e}, M\"{u}ller, Fran\c{c}ois,
\AA{}ngstr\"{o}m, Erd\H{o}s, \O{}rsted, stra\ss{}e, na\"{\i}ve. And the same
words typed as literal UTF-8: René, Müller, François, Ångström.

Text-mode super/subscripts: the 3\textsuperscript{rd} entry of
H\textsubscript{2}O, and inline verbatim like \verb|x <- y| or
\verb+a & b+ that must not be interpreted.

\section{Mathematics}
\label{sec:math}

Inline math $\norm{x}_2 \le \sqrt{d}\,\norm{x}_\infty$ sits in running text,
as does $\Tr(A) = \sum_i A_{ii}$ and the operator $\argmin_{x \in \R^d} f(x)$.

A numbered equation with a label:
\begin{equation}
  \label{eq:objective}
  f(x) = \tfrac{1}{2}\norm{Ax - b}^2 + \lambda \norm{x}_1 .
\end{equation}

Grouped equations that share a number, each line separately labelled:
\begin{subequations}
  \label{eq:system}
  \begin{align}
    \dot{x} &= \sigma (y - x),          \label{eq:lorenz-x} \\
    \dot{y} &= x(\rho - z) - y,         \label{eq:lorenz-y} \\
    \dot{z} &= xy - \beta z .           \label{eq:lorenz-z}
  \end{align}
\end{subequations}

A case analysis and a matrix, both inside a display:
\begin{equation}
  \label{eq:cases}
  \operatorname{sgn}(t) =
  \begin{cases}
    +1 & t > 0, \\
     0 & t = 0, \\
    -1 & t < 0,
  \end{cases}
  \qquad
  M = \begin{pmatrix}
    1 & 0 \\
    0 & 1
  \end{pmatrix}.
\end{equation}

Alignment across three columns with \texttt{alignat}, plus a \texttt{gather}:
\begin{alignat}{2}
  a &= b + c    &\quad& \text{(definition)} \label{eq:alignat-1} \\
  d &= e - f    &&      \text{(consequence)} \label{eq:alignat-2}
\end{alignat}

\begin{gather*}
  p + q = r \\
  s + t = u
\end{gather*}

An unnumbered display with \verb|\[ ... \]|: \[ \int_0^\infty e^{-x^2}\,dx =
\frac{\sqrt{\pi}}{2}. \]

By \cref{eq:objective,eq:cases} and \eqref{eq:lorenz-y} the claim follows; see
also \autoref{sec:math} and page~\pageref{sec:tables}.

\section{Tables}
\label{sec:tables}

\begin{table}[htbp]
  \centering
  \caption[Short caption for the list of tables]{Booktabs table with a
  multi-column header and a multi-row stub.}
  \label{tab:results}
  \begin{tabular}{@{}llrr@{}}
    \toprule
    \multirow{2}{*}{Method} & \multicolumn{3}{c}{Benchmark} \\
    \cmidrule(l){2-4}
                            & Split & Error (\%) & Time (s) \\
    \midrule
    Baseline                & dev   & 12.4       & 0.81 \\
                            & test  & 13.1       & 0.79 \\
    Ours                    & dev   & \textbf{9.2} & 1.04 \\
                            & test  & \textbf{9.8} & 1.02 \\
    \bottomrule
  \end{tabular}
\end{table}

A table with vertical rules and an inline math cell:

\begin{table}[t]
  \centering
  \begin{tabular}{|l|c|}
    \hline
    Symbol & Meaning \\
    \hline
    $\R^d$      & Euclidean space \\
    $\norm{\cdot}$ & A norm on $\R^d$ \\
    \hline
  \end{tabular}
  \caption{Notation.}
  \label{tab:notation}
\end{table}

\section{Figures}
\label{sec:figures}

\begin{figure}[t]
  \centering
  \includegraphics[width=0.6\linewidth]{plots/convergence.pdf}
  \caption{A single-image figure.}
  \label{fig:convergence}
\end{figure}

\begin{figure}[t]
  \centering
  \begin{subfigure}[b]{0.45\linewidth}
    \centering
    \includegraphics[width=\linewidth]{plots/left.pdf}
    \caption{Left panel.}
    \label{fig:left}
  \end{subfigure}
  \hfill
  \begin{subfigure}[b]{0.45\linewidth}
    \centering
    \includegraphics[width=\linewidth]{plots/right.pdf}
    \caption{Right panel.}
    \label{fig:right}
  \end{subfigure}
  \caption{Two panels side by side.}
  \label{fig:panels}
\end{figure}

\section{Lists}
\label{sec:lists}

Nested lists have to keep their nesting:

\begin{enumerate}[label=(\roman*)]
  \item A first outer item.
  \item A second outer item with an inner list:
    \begin{itemize}
      \item An inner bullet.
      \item Another inner bullet, with \emph{emphasis} and $x^2$.
    \end{itemize}
  \item A third outer item.
\end{enumerate}

A description list:

\begin{description}
  \item[Convex] The epigraph is a convex set.
  \item[Smooth] The gradient is Lipschitz.
\end{description}

A multi-paragraph list item:

\begin{itemize}
  \item First paragraph of the item.

    Second paragraph of the same item.
  \item A short item.
\end{itemize}

\section{Theorems and Proofs}
\label{sec:theorems}

\begin{definition}[Strong convexity]
  \label{def:strong}
  A function $f$ is $\mu$-strongly convex if $f - \tfrac{\mu}{2}\norm{\cdot}^2$
  is convex.
\end{definition}

\begin{theorem}[Convergence rate]
  \label{thm:rate}
  Let $f$ satisfy \cref{def:strong}. Then
  \begin{equation}
    \label{eq:rate}
    f(x_T) - f^\star \le \frac{C}{T} .
  \end{equation}
\end{theorem}

\begin{proof}
  Combine \eqref{eq:objective} with the descent lemma. The remaining steps are
  standard \citep{boyd2004convex}.
\end{proof}

\section{Verbatim and Listings}
\label{sec:code}

\begin{verbatim}
def f(x):
    # a comment with $math$ and \backslash
    return x ** 2
\end{verbatim}

\begin{lstlisting}[language=Python,caption={A listing}]
for i in range(10):
    print(i)   # 100% literal
\end{lstlisting}

\section{Citations and Notes}
\label{sec:cites}

Standard forms: \cite{knuth1984texbook}, \citep{boyd2004convex},
\citet{nesterov2018lectures}, and a page-qualified
\citep[see][p.~42]{boyd2004convex}. Multiple keys collapse:
\citep{boyd2004convex,nesterov2018lectures}.

A footnote with structure inside it.\footnote{Footnotes may contain math like
$e^{i\pi} + 1 = 0$ and citations~\citep{knuth1984texbook}.}

\appendix

\section{Deferred Proofs}
\label{app:proofs}

\begin{lemma}
  \label{lem:aux}
  An auxiliary bound.
\end{lemma}

The appendix section is lettered, so \cref{app:proofs} renders as ``Appendix A''.

\begin{thebibliography}{9}

\bibitem{boyd2004convex}
S.~Boyd and L.~Vandenberghe.
\newblock \emph{Convex Optimization}.
\newblock Cambridge University Press, 2004.

\bibitem{knuth1984texbook}
D.~E. Knuth.
\newblock \emph{The \TeX{}book}.
\newblock Addison-Wesley, 1984.

\bibitem{nesterov2018lectures}
Y.~Nesterov.
\newblock \emph{Lectures on Convex Optimization}.
\newblock Springer, 2018.

\end{thebibliography}

\end{document}
